\chapter{Orders}\label{ch:orders}

\begin{center}\itshape
Every sequence of $n^2+1$ distinct numbers contains a monotone subsequence of length $n+1$.\\
\upshape(Erd\H os and Szekeres, 1935, in our words)
\end{center}

Budapest, around 1933. A group of young people meets regularly to trade
problems; among them are Paul Erd\H os, about twenty, George Szekeres, a student of chemical engineering, and Esther Klein.
Klein brings them a fact she has found: among five points of the plane, no three on a line, four always form a
convex quadrilateral. She explains why. Then the question that would not let go: for a
convex pentagon, how many points are needed? For a hexagon? Szekeres found that the answer exists, by rediscovering
Ramsey's theorem on the way, and Erd\H os found a second proof, built on a lemma about sequences: the epigraph of this
chapter. The paper appeared in 1935. Szekeres and Klein married in 1937, and Erd\H os, who loved the story, called the
whole subject \emph{the happy ending problem}.

The lemma on sequences is a theorem about \emph{orders}: two orders on the same terms, by position and by size,
that must agree for a long stretch or disagree for a long stretch. Orders are everywhere in computing. A build system
compiles files in an order compatible with their dependencies; a scheduler runs tasks after their prerequisites; a
compiler computes the meaning of a recursive definition as the least solution of an equation; a proof that a program
terminates exhibits a quantity that decreases in a well-ordered set. This chapter develops the mathematics behind all
of these: partial orders and lattices, the fixed point theorem of Knaster and Tarski, well-founded orders and the art
of proving termination, the theorems of Mirsky and Dilworth on chains and antichains, and the beautiful structure of
the lattice of all subsets of a set.

\begin{think}
Write down any ten distinct numbers in a row. Can you always find four of them, read from left to right, that increase,
or four that decrease? Try to build a row of nine numbers without such a quadruple. Then try ten.
\end{think}

Two earlier chapters already met a monotone piece of a sequence. Chapter~\ref{ch:counting} proved the epigraph by the
pigeonhole principle, labelling each term by the lengths of the longest increasing and decreasing subsequences that end
there. Chapter~\ref{ch:sequences} builds one monotone subsequence from the peaks of a bounded sequence, so that completeness
can pass to a limit. This chapter asks what an order is. The epigraph returns as a theorem about chains and antichains, and
the same language proves Dilworth's theorem and Sperner's theorem.

\section{Relations and partial orders}

A \emph{relation} on a set $P$ is a subset $R\subseteq P\times P$; we write $x\,R\,y$ for $(x,y)\in R$. An \emph{equivalence relation} is reflexive
($x\,R\,x$), symmetric ($x\,R\,y\Rightarrow y\,R\,x$) and transitive ($x\,R\,y\,R\,z\Rightarrow x\,R\,z$). Its classes form a partition of $P$, and every partition
comes from exactly one equivalence relation, ``being in the same block''.

\begin{definition}\label{orders:def:poset}
A \emph{partial order} is a relation $\le$ that is reflexive, antisymmetric ($x\le y\le x\Rightarrow x=y$) and transitive; $(P,\le)$ is a
\emph{poset}. We write $x<y$ for $x\le y$, $x\ne y$. Two elements are \emph{comparable} if $x\le y$ or $y\le x$. A \emph{chain} is a set of pairwise
comparable elements, an \emph{antichain} a set of pairwise incomparable ones. The order is \emph{total} (or linear) if any two
elements are comparable. An element $m$ is \emph{minimal} if no $x<m$, and \emph{least} if $m\le x$ for all $x$; similarly maximal and
greatest. We say $y$ \emph{covers} $x$ if $x<y$ and no $z$ satisfies $x<z<y$.
\end{definition}

Examples: the integers with divisibility; the subsets of a set with inclusion; the tasks of a project with ``must come
before''; the vectors of $\R^n$ with the coordinatewise order; words with the dictionary order (total) or the prefix order
(not total). A finite poset is drawn as a \emph{Hasse diagram}: one point per element, higher points for larger elements, and
a segment for each covering pair. Since $x<y$ holds exactly when a path of coverings climbs from $x$ to $y$ (in a finite poset,
refine $x<y$ by inserting elements until no more fit), the diagram determines the order.

\begin{figure}[ht]
\centering
\fig{orders-ch02-divisors}
\caption{The divisors of $60$ ordered by divisibility. Each segment goes up by multiplying by one prime. The lattice is the
product of three chains, $\{1,2,4\}$, $\{1,3\}$ and $\{1,5\}$, and its middle rank (shaded), the divisors with two prime factors $\{4,6,10,15\}$, is an antichain of maximal size
(Exercise~\ref{orders:ex:divisorscd}).}
\label{orders:fig:divisors}
\end{figure}

A \emph{linear extension} of a poset is a total order $\preceq$ on the same set with $x\le y\Rightarrow x\preceq y$: a way of doing all the tasks one
after the other while respecting the prerequisites.

\begin{proposition}[topological sorting]\label{orders:prop:topsort}
Every finite poset has a linear extension.
\end{proposition}
\begin{proof}
A finite non-empty poset has a minimal element: start anywhere and go down while possible; as the elements met are
strictly decreasing they are distinct, so this stops. Take a minimal element $m_1$ of $P$, then a minimal element $m_2$ of
$P\setminus\{m_1\}$, and so on, and order $m_1\prec m_2\prec\dots\prec m_n$. If $m_i<m_j$ in $P$, then $j>i$: when $m_j$ was chosen it was minimal among the
remaining elements, so $m_i$ was no longer there.
\end{proof}

This is the algorithm run by \texttt{make} and by every package manager. For infinite
posets, every partial order still extends to a total one (Szpilrajn, 1930), by Zorn's lemma (Volume~6, Chapter~1).

\section{Lattices and fixed points}

\begin{definition}\label{orders:def:lattice}
A poset is a \emph{lattice} if any two elements $x,y$ have a least upper bound $x\lor y$ (their \emph{join}) and a greatest lower bound
$x\land y$ (their \emph{meet}). It is a \emph{complete lattice} if every subset $S$, even infinite or empty, has a least upper bound $\bigvee S$
and a greatest lower bound $\bigwedge S$.
\end{definition}

The subsets of a set form a complete lattice ($\bigcup$ and $\bigcap$); so do the divisors of $n$ (lcm and gcd), the subspaces of a
vector space (sum and intersection), the partitions of a set ordered by refinement, and the interval $[0,1]$. A finite
lattice is complete: take iterated joins and meets, with $\bigvee\emptyset$ the least element (the meet of everything) and $\bigwedge\emptyset$ the
greatest. Boolean algebras are the lattices with complements and distributivity.

In 1928 Bronis\l aw Knaster, in Warsaw, proved a small lemma about sets; in 1955 Alfred Tarski stated it for lattices.
It has since become the foundation of the semantics of programming languages.

\begin{theorem}[Knaster and Tarski]\label{orders:thm:tarski}
Let $L$ be a complete lattice and $f\colon L\to L$ monotone ($x\le y\Rightarrow f(x)\le f(y)$). Then $f$ has a least fixed point, namely
$a=\bigwedge\{x:f(x)\le x\}$, and a greatest fixed point. More generally, the fixed points of $f$ form a complete lattice.
\end{theorem}
\begin{proof}
Let $D=\{x:f(x)\le x\}$, which contains the greatest element, and $a=\bigwedge D$. For $x\in D$, $a\le x$ so $f(a)\le f(x)\le x$; thus $f(a)$ is a lower bound
of $D$ and $f(a)\le a$. Then $f(f(a))\le f(a)$, so $f(a)\in D$ and $a\le f(a)$. Hence $f(a)=a$. Every fixed point lies in $D$, so it is $\ge a$: $a$ is the
least fixed point. The greatest one is obtained dually, as $\bigvee\{x:x\le f(x)\}$.

Now let $S$ be a set of fixed points and $s=\bigvee S$. For $y\in S$, $y=f(y)\le f(s)$, so $s\le f(s)$, and $f$ maps the interval $[s,\top]=\{x:x\ge s\}$
into itself (if $x\ge s$ then $f(x)\ge f(s)\ge s$). This interval is a complete lattice, so $f$ has a least fixed point $b$ in it. Every
fixed point that is an upper bound of $S$ lies in $[s,\top]$ and so is $\ge b$: $b$ is the least upper bound of $S$ among fixed points.
Greatest lower bounds are obtained dually.
\end{proof}

The fixed points need not form a \emph{sublattice}: the join of two fixed points among fixed points may be larger than their
join in $L$. A famous consequence is the Schr\"oder--Bernstein theorem, which Cantor had stated without being able to prove.

\begin{corollary}[Schr\"oder--Bernstein]\label{orders:cor:sb}
If there are injections $f\colon A\to B$ and $g\colon B\to A$, there is a bijection between $A$ and $B$.
\end{corollary}
\begin{proof}
On the complete lattice of subsets of $A$, the map $\Phi(S)=A\setminus g(B\setminus f(S))$ is monotone: a larger $S$ makes $f(S)$ larger, $B\setminus f(S)$
smaller, its image under $g$ smaller, and the complement larger. Let $S$ be a fixed point, so that $A\setminus S=g(B\setminus f(S))$. Define $h(x)=f(x)$ for
$x\in S$ and $h(x)=g^{-1}(x)$ for $x\in A\setminus S$, which makes sense since every such $x$ is $g(y)$ for exactly one $y\in B\setminus f(S)$. Then $h$ maps $S$
bijectively onto $f(S)$ and $A\setminus S$ bijectively onto $B\setminus f(S)$; these two images partition $B$, so $h$ is a bijection.
\end{proof}

On a finite lattice the least fixed point can be computed by iteration, and this is how compilers and databases do
it.

\begin{proposition}[Kleene iteration]\label{orders:prop:kleene}
Let $L$ be a finite lattice with least element $\bot$ and $f$ monotone. The sequence $\bot\le f(\bot)\le f^2(\bot)\le\cdots$ is non-decreasing and
becomes constant, equal to the least fixed point of $f$, after at most $h$ steps, where $h+1$ is the largest size of a chain in $L$.
\end{proposition}
\begin{proof}
$\bot\le f(\bot)$, and applying $f$ repeatedly gives $f^k(\bot)\le f^{k+1}(\bot)$ by monotonicity. A strictly increasing stretch of this sequence
is a chain, so after at most $h$ strict increases the sequence stops increasing, and then $f^k(\bot)=f^{k+1}(\bot)$ is a fixed point. If $x$ is
any fixed point, $\bot\le x$ gives by induction $f^k(\bot)\le f^k(x)=x$; so the limit is the least fixed point.
\end{proof}

\begin{example}\label{orders:ex:reach}
\emph{Reachability.} In a directed graph with vertex set $V$ and a start vertex $s$, consider on the lattice of subsets of $V$ the
monotone map $f(X)=\{s\}\cup\{v:\text{some }u\in X\text{ has an arc }u\to v\}$. Its least fixed point is the set of vertices reachable from $s$, and Kleene
iteration is breadth-first search: $f^k(\emptyset)$ is the set of vertices at distance less than $k$. The least model of a set of Horn
clauses is the least fixed point of ``add every conclusion whose hypotheses are true'', and recursive
queries in SQL are evaluated in the same way.
\end{example}

\section{Well-founded orders and termination}

Why does Euclid's algorithm stop? Because the second argument is a natural number that decreases at each step, and
natural numbers cannot decrease forever. Every termination proof has this shape; only the set where the quantity
decreases changes.

\begin{definition}\label{orders:def:wf}
A poset is \emph{well-founded} if every non-empty subset has a minimal element. A total well-founded order is a
\emph{well-order}.
\end{definition}

A well-founded poset has no infinite strictly decreasing sequence $x_0>x_1>\cdots$, since the set of its terms would have no
minimal element. Conversely, if a non-empty subset $S$ has no minimal element, we can choose $x_0\in S$, then $x_1<x_0$ in $S$, and so
on, which gives an infinite decreasing sequence; this step-by-step choice is the axiom of dependent choice (Volume~6,
Chapter~1), and we use the two forms of the definition interchangeably. For instance $\N$, finite posets and $\N^k$ with
the coordinatewise order are well-founded; $\Z$ and $\Q_{\ge0}$ are not.

\begin{theorem}[well-founded induction]\label{orders:thm:wfind}
Let $(P,<)$ be well-founded and $\mathcal P(x)$ a property of the elements of $P$. If, for every $x$, $\mathcal P(y)$ for all $y<x$ implies $\mathcal P(x)$,
then $\mathcal P(x)$ holds for all $x\in P$.
\end{theorem}
\begin{proof}
Otherwise the set of counterexamples has a minimal element $x$; then $\mathcal P(y)$ holds for all $y<x$, so $\mathcal P(x)$ holds, a
contradiction.
\end{proof}

Ordinary induction is the case $P=\N$. To build well-founded orders from simpler ones, the \emph{lexicographic order} on $A\times B$
says $(a,b)<(a',b')$ if $a<a'$, or $a=a'$ and $b<b'$.

\begin{proposition}\label{orders:prop:lex}
If $A$ and $B$ are well-founded, so is $A\times B$ with the lexicographic order.
\end{proposition}
\begin{proof}
Let $S\subseteq A\times B$ be non-empty. Let $a_0$ be a minimal element of the set of first coordinates of $S$, and $b_0$ a minimal element of
$\{b:(a_0,b)\in S\}$. If $(a,b)<(a_0,b_0)$ with $(a,b)\in S$, then either $a<a_0$, against the choice of $a_0$, or $a=a_0$ and $b<b_0$, against the choice of $b_0$.
So $(a_0,b_0)$ is minimal in $S$.
\end{proof}

\begin{example}\label{orders:ex:ackermann}
\emph{Ackermann's function} is defined by $A(0,n)=n+1$, $A(m+1,0)=A(m,1)$ and $A(m+1,n+1)=A(m,A(m+1,n))$. Does the recursion
always stop? Each recursive call is made on a pair that is lexicographically smaller: $(m,1)<(m+1,0)$, $(m+1,n)<(m+1,n+1)$,
and $(m,k)<(m+1,n+1)$ whatever $k$ is. By well-founded induction on $\N\times\N$, $A(m,n)$ is defined for all $m,n$. Its growth is
spectacular: $A(2,n)=2n+3$, $A(3,n)=2^{n+3}-3$, and $A(4,2)=2^{65536}-3$ has $19\,729$ digits (Exercise~\ref{orders:ex:ackermannvalues}). Its slowly
growing inverse governs the union--find data structure used for disjoint-set maintenance.
\end{example}

Sometimes no single quantity decreases, but a \emph{bag} of quantities does: one element is removed and replaced by any
number of smaller ones. Think of a to-do list in which finishing a task may create many new, smaller tasks. Will the list
ever be empty? A \emph{multiset} over $X$ is a finite collection of elements of $X$ with repetitions allowed. For multisets $M,N$, write
$N<_{\mathrm{mul}}M$ if $N$ is obtained from $M$ by removing a non-empty submultiset $R$ and adding a multiset $Q$ each of whose elements is
smaller than some element of $R$.

\begin{theorem}[Dershowitz and Manna, 1979]\label{orders:thm:multiset}
If $(X,<)$ has no infinite decreasing sequence, then neither has the multiset order: every sequence of multisets
$M_0>_{\mathrm{mul}}M_1>_{\mathrm{mul}}\cdots$ is finite.
\end{theorem}
\begin{proof}
Suppose $M_0>_{\mathrm{mul}}M_1>_{\mathrm{mul}}\cdots$ is infinite. Grow a rooted tree whose vertices, apart from the root, are the individual copies of
elements that appear: the copies in $M_0$ are the children of the root, and when step $i$ removes $R_i$ and adds $Q_i$, attach every
copy in $Q_i$ as a child of one copy in $R_i$ that is larger than it. Each vertex gets children at most once, when it is
removed, and finitely many; so every vertex has finitely many children. Each step removes at least one copy, never
the same one twice, so the tree is infinite.

By K\H onig's lemma (Exercise~\ref{orders:ex:konig}, whose direct proof is two lines: always move to a child with infinitely many
descendants), the tree has an infinite path from the root, $x_1,x_2,\dots$, and each child is smaller than its parent:
$x_1>x_2>\cdots$ is an infinite decreasing sequence in $X$, a contradiction.
\end{proof}

The to-do list empties. The theorem is used inside proof assistants and rewriting systems to show that programs and
simplification rules terminate. Exercise~\ref{orders:ex:hydra} uses it to kill the Hydra of Kirby and Paris, a monster whose
heads grow back by the hundreds and which, nevertheless, always dies.

\section{Chains and antichains}

Back to Budapest. In a poset, a chain and an antichain meet in at most one point, so if $P$ is covered by $k$ chains it has
no antichain larger than $k$, and if it is covered by $k$ antichains it has no chain larger than $k$. The two theorems of this
section say that in both cases the obvious obstruction is the only one. The \emph{height} of a finite poset is the size of
its largest chain, its \emph{width} the size of its largest antichain.

\begin{theorem}[Mirsky, 1971]\label{orders:thm:mirsky}
A finite poset of height $h$ can be partitioned into $h$ antichains, and not into fewer.
\end{theorem}
\begin{proof}
Fewer is impossible, since a chain of size $h$ meets each antichain at most once. For $x\in P$ let $r(x)$ be the size of the largest
chain whose greatest element is $x$; then $1\le r(x)\le h$. If $x<y$, a largest chain ending at $x$ can be continued by $y$, so $r(y)>r(x)$. Hence
each level set $\{x:r(x)=i\}$ is an antichain, and the $h$ level sets partition $P$.
\end{proof}

The dual statement is harder, and Robert Dilworth's proof of 1950 made it a classic.

\begin{theorem}[Dilworth, 1950]\label{orders:thm:dilworth}
A finite poset of width $w$ can be partitioned into $w$ chains, and not into fewer.
\end{theorem}
\begin{proof}[Proof (Perles, 1963)]
Fewer is impossible, since an antichain of size $w$ meets each chain at most once. We show by induction on $|P|$ that $w$
chains suffice; the empty poset is clear. Let $a$ be a minimal element of $P$ and $b$ a maximal element with $a\le b$ (go up from $a$
while possible). The set $C=\{a,b\}$ is a chain. If $P\setminus C$ has width at most $w-1$, it is a union of $w-1$ chains by induction, and $C$
completes the partition. Otherwise $P\setminus C$ contains an antichain $A$ of size $w$. Let
\[
\begin{gathered}P^+=\{x\in P:x\ge y\text{ for some }y\in A\},\qq P^-=\{x\in P:x\le y\text{ for some }y\in A\}.\end{gathered}
\]
Every $x\in P$ is comparable with some element of $A$, otherwise $A\cup\{x\}$ would be an antichain of size $w+1$; so $P^+\cup P^-=P$. If $x\in P^+\cap P^-$,
then $y\le x\le y'$ with $y,y'\in A$, so $y=y'=x$: thus $P^+\cap P^-=A$. Moreover $a\notin P^+$: $a\ge y\in A$ would force $a=y$ by minimality, but $a\notin A$.
Likewise $b\notin P^-$. So $P^+$ and $P^-$ are smaller than $P$, both have width $w$ (they contain $A$), and by induction each is a union of $w$
chains. Each of the $w$ chains of $P^+$ contains exactly one element of $A$, and that element is its least element: if $y\in A$ and $z<y$
were in the same chain of $P^+$, then $z\ge y'$ for some $y'\in A$, and $y'<y$ contradicts that $A$ is an antichain. Likewise each chain of
$P^-$ has exactly one element of $A$, as its greatest element. Gluing the chain of $P^-$ ending at $y$ to the chain of $P^+$ starting at
$y$, for each $y\in A$, gives $w$ chains covering $P$.
\end{proof}

A second proof uses K\H onig's theorem on matchings. Now
the happy ending.

\begin{theorem}[Erd\H os and Szekeres, 1935]\label{orders:thm:es}
Every sequence of $mn+1$ distinct real numbers contains an increasing subsequence of length $m+1$ or a decreasing
subsequence of length $n+1$. The bound $mn+1$ is sharp.
\end{theorem}
\begin{proof}
Let the sequence be $a_1,\dots,a_N$ with $N=mn+1$, and order the indices by $i\preceq j$ if $i\le j$ and $a_i\le a_j$. This is a partial order whose
chains are the increasing subsequences and whose antichains are the decreasing ones (if $i<j$ are incomparable, then
$a_i>a_j$). If there is no increasing subsequence of length $m+1$, the height is at most $m$, and Mirsky's theorem covers the $mn+1$
indices by $m$ antichains; one of them has at least $n+1$ elements, a decreasing subsequence.

For sharpness, take $m$ decreasing blocks of $n$ numbers, each block above the previous one, for instance for $m=n=3$:
$3,2,1,\;6,5,4,\;9,8,7$. An increasing subsequence takes at most one term from each block, and a decreasing one stays in one
block.
\end{proof}

\begin{figure}[ht]
\centering
\fig{orders-ch02-es}
\caption{A sequence of $10$ distinct numbers (charcoal dots, position across, value up). Each term is labelled with the
length of the longest increasing subsequence ending at it, the function $r$ of Mirsky's proof. Three terms share
the label $2$, and they decrease from left to right: equal labels always form a decreasing subsequence.}
\label{orders:fig:es}
\end{figure}

\begin{remark}
The \emph{happy ending problem} itself asks for the least $f(n)$ such that any $f(n)$ points in the plane, no three on a line,
contain the vertices of a convex $n$-gon. Klein's observation is $f(4)=5$ (Exercise~\ref{orders:ex:klein}); $f(5)=9$ was proved by Endre Makai, and $f(6)=17$ by a computer search of Szekeres and Peters in 2006, the year after Szekeres and Klein died within an
hour of each other. Erd\H os and Szekeres conjectured $f(n)=2^{n-2}+1$, and proved $f(n)\le\binom{2n-4}{n-2}+1$. In 2017 Andrew Suk proved
$f(n)=2^{n+o(n)}$, a near-miss of the conjecture that is still open.
\end{remark}

\section{The Boolean lattice}

The subsets of $\{1,\dots,n\}$, ordered by inclusion, form the \emph{Boolean lattice} $B_n$; its elements of rank $k$ are the $k$-subsets. How
many subsets can be chosen so that none contains another? Taking all the subsets of size $\lfloor n/2\rfloor$ gives $\binom n{\lfloor n/2\rfloor}$. Emanuel
Sperner, a student in Hamburg, proved in 1928 that one cannot do better. Here is the proof that Lubell, Yamamoto and
Meshalkin found independently, a counting argument that weighs each set by the number of ways to reach it.

\begin{theorem}[LYM inequality]\label{orders:thm:lym}
If $\mathcal F$ is an antichain in $B_n$, then $\displaystyle\sum_{A\in\mathcal F}\binom n{|A|}^{-1}\le1$. Consequently $|\mathcal F|\le\binom n{\lfloor n/2\rfloor}$ (Sperner's theorem).
\end{theorem}
\begin{proof}
A \emph{maximal chain} $\emptyset\subset\{\sigma_1\}\subset\{\sigma_1,\sigma_2\}\subset\dots\subset\{1,\dots,n\}$ is determined by a permutation $\sigma$ of $\{1,\dots,n\}$, so there are $n!$ of them. A set $A$
of size $k$ lies on the maximal chains whose permutation lists the elements of $A$ first, in any order, and then the others:
there are $k!(n-k)!$ of them. Since a chain meets the antichain $\mathcal F$ at most once, these sets of chains are disjoint for $A\in\mathcal F$, so
$\sum_{A\in\mathcal F}|A|!\,(n-|A|)!\le n!$. Dividing by $n!$ gives the inequality. Since every binomial coefficient $\binom nk$ is at most $\binom n{\lfloor n/2\rfloor}$, each term of
the sum is at least $\binom n{\lfloor n/2\rfloor}^{-1}$, and $|\mathcal F|\le\binom n{\lfloor n/2\rfloor}$.
\end{proof}

By Dilworth's theorem, Sperner's theorem says that $B_n$ is a union of $\binom n{\lfloor n/2\rfloor}$ chains. These chains can be chosen
remarkably well. A chain $A_1\subset A_2\subset\dots\subset A_m$ in $B_n$ is \emph{saturated} if $|A_{i+1}|=|A_i|+1$, and \emph{symmetric} if moreover $|A_1|+|A_m|=n$: it climbs
from rank $k$ to rank $n-k$ without skipping a rank.

\begin{theorem}[de Bruijn, van Ebbenhorst Tengbergen and Kruyswijk, 1951]\label{orders:thm:scd}
$B_n$ is the disjoint union of symmetric saturated chains; there are $\binom n{\lfloor n/2\rfloor}$ of them.
\end{theorem}
\begin{proof}
Induction on $n$; for $n=1$, the single chain $\emptyset\subset\{1\}$ works. Given such a decomposition of $B_{n-1}$, replace each of its chains
$A_1\subset\dots\subset A_m$ by two chains of $B_n$:
\[
\begin{gathered}A_1\subset\dots\subset A_m\subset A_m\cup\{n\}\qq\text{and, if }m\ge2,\qq A_1\cup\{n\}\subset\dots\subset A_{m-1}\cup\{n\}.\end{gathered}
\]
If the old chain ran from rank $k$ to rank $n-1-k$, the first new chain runs from $k$ to $n-k$ and the second from $k+1$ to $n-1-k$; both are
saturated and symmetric in $B_n$. Every subset of $\{1,\dots,n\}$ appears exactly once: a set $A$ not containing $n$ lies in exactly one old chain,
hence in the first new chain built from it; a set $A\cup\{n\}$ lies in the first new chain if $A$ is the top of its old chain, and in the
second otherwise. Finally every symmetric saturated chain meets the middle rank $\lfloor n/2\rfloor$ exactly once, so the number of chains is
$\binom n{\lfloor n/2\rfloor}$.
\end{proof}

\begin{figure}[ht]
\centering
\fig{orders-ch02-scd}
\caption{A symmetric chain decomposition of $B_4$ built by the proof of Theorem~\ref{orders:thm:scd}: six chains (thick), one from $\emptyset$ to
$\{1,2,3,4\}$, three from rank $1$ to rank $3$, two inside the middle rank. The thin segments are the other covering
relations. The same construction extends to $B_5$.}
\label{orders:fig:scd}
\end{figure}

\begin{beauty}
One recursive rule, ``each chain splits into a longer one and a shorter one'', organises all $2^n$ subsets into
$\binom n{\lfloor n/2\rfloor}$ ladders standing symmetrically about the middle. It proves Sperner's theorem with no computation, shows that the
binomial coefficients are unimodal ($\binom nk\le\binom n{k+1}$ for $k<n/2$, since each chain through rank $k$ continues to rank $k+1$), and gives a
fast algorithm, which Exercise~\ref{orders:ex:hansel} turns into a bound on the number of monotone Boolean functions.
\end{beauty}

One last classic. A family of sets is \emph{intersecting} if any two of its members meet. All $k$-subsets of $\{1,\dots,n\}$ containing
the element $1$ form an intersecting family of size $\binom{n-1}{k-1}$. Paul Erd\H os, Chao Ko and Richard Rado proved in 1938 that one
cannot do better when $n\ge2k$, but did not publish until 1961, as nobody seemed to care about such questions; today it is
one of the most cited theorems of combinatorics. The proof below, by Gyula Katona (1972), puts the elements around a
round table.

\begin{theorem}[Erd\H os, Ko and Rado]\label{orders:thm:ekr}
If $n\ge2k$ and $\mathcal F$ is an intersecting family of $k$-subsets of $\{1,\dots,n\}$, then $|\mathcal F|\le\binom{n-1}{k-1}$.
\end{theorem}
\begin{proof}
Seat $1,\dots,n$ around a round table in some cyclic order $\sigma$; an \emph{arc} is a set of $k$ consecutive seats. We claim that at most $k$ arcs
of $\sigma$ belong to $\mathcal F$. Suppose $A\in\mathcal F$ is an arc, occupying seats $1,\dots,k$. Every other arc meeting $A$ either starts at a seat
$i+1$ or ends at a seat $i$, for some $1\le i\le k-1$, and for each such $i$ these two arcs are disjoint, since together they occupy
$2k\le n$ consecutive seats. As $\mathcal F$ is intersecting, it contains at most one of the two, so at most $1+(k-1)=k$ arcs in all.

Now count pairs $(\sigma,A)$ with $\sigma$ a cyclic order and $A\in\mathcal F$ an arc of $\sigma$. There are $(n-1)!$ cyclic orders, so at most $k\,(n-1)!$ pairs.
A given $k$-set $A$ is an arc of exactly $k!\,(n-k)!$ cyclic orders: seat the elements of $A$ in a block in one of $k!$ orders, then the
others in $(n-k)!$ orders, and count cyclic orders by fixing the block's position. So $|\mathcal F|\,k!\,(n-k)!\le k\,(n-1)!$, that is,
$|\mathcal F|\le\frac{k\,(n-1)!}{k!\,(n-k)!}=\binom{n-1}{k-1}$.
\end{proof}

For $n<2k$ the bound fails, as any two $k$-subsets then meet and all $\binom nk$ of them form an intersecting family.

\section{Exercises}

\begin{exercise}[$\star$]\label{orders:ex:klein}
\emph{Esther Klein's observation.} Show that among any five points of the plane, no three on a line, four are the vertices of
a convex quadrilateral.
\end{exercise}

\begin{exercise}[$\star$]\label{orders:ex:esproduct}
Show that in every sequence of $N$ distinct numbers, the length of the longest increasing subsequence times the length of
the longest decreasing subsequence is at least $N$.
\end{exercise}

\begin{exercise}[$\star$]\label{orders:ex:ackermannvalues}
Show that $A(1,n)=n+2$, $A(2,n)=2n+3$ and $A(3,n)=2^{n+3}-3$, and that $A(4,n)+3$ is a tower of $n+3$ twos. How many decimal digits has
$A(4,2)$?
\end{exercise}

\begin{exercise}[$\star\star$]\label{orders:ex:rewriting}
Start with a word over the letters $a,b$ and repeatedly replace some occurrence of $ba$ by $aab$. Show that this always
stops, whatever the choices, and that from $b^ka$ it stops at $a^{2^k}b^k$ after exactly $2^k-1$ steps.
\end{exercise}

\begin{exercise}[$\star\star\star$]\label{orders:ex:hydra}
\emph{Hercules and the Hydra} (Kirby and Paris, 1982). A hydra is a finite rooted tree; its heads are the leaves other than the
root. At step $n$, Hercules cuts a head $h$; let $p$ be its parent. If $p$ is the root, nothing grows. Otherwise, let $g$ be the parent
of $p$: the hydra grows $n$ new copies of the subtree at $p$ (as it is after the cut), attached to $g$. Show that Hercules wins, that is,
the hydra is reduced to its root after finitely many steps, whatever heads he cuts.
\end{exercise}

\begin{exercise}[$\star$]\label{orders:ex:fixedpoints}
(a) Give a map of the two-element lattice $\{0,1\}$ to itself with no fixed point. (b) Give a monotone map $f$ of $[0,1]$ to itself
whose least fixed point is not reached by the iterates $f^k(0)$ after finitely many steps. Is it their limit?
\end{exercise}

\begin{exercise}[$\star\star$]\label{orders:ex:youngtableaux}
Show that the number of linear extensions of the poset $\{1,2\}\times\{1,\dots,n\}$ with the coordinatewise order is the Catalan number
$\frac1{n+1}\binom{2n}n$.
\end{exercise}

\begin{exercise}[$\star$]\label{orders:ex:dedekind}
Show that the antichains of $B_n$ are in bijection with the monotone Boolean functions of $n$ variables. Count them for $n=2$
and $n=3$.
\end{exercise}

\begin{exercise}[$\star\star$]\label{orders:ex:hansel}
Let $D(n)$ be the number of antichains of $B_n$. Show that $2^{\binom n{\lfloor n/2\rfloor}}\le D(n)\le(n+2)^{\binom n{\lfloor n/2\rfloor}}$, so that $\log_2D(n)$ is $\binom n{\lfloor n/2\rfloor}$ up to a factor
$\log_2(n+2)$.
\end{exercise}

\begin{exercise}[$\star\star\star$]\label{orders:ex:divisorscd}
Let $n=p_1^{a_1}\cdots p_r^{a_r}$. Show that the divisors of $n$, ordered by divisibility, can be partitioned into symmetric
saturated chains (rank being the number of prime factors counted with multiplicity), and deduce that the largest
antichain of divisors of $n$ consists of the divisors with $\lfloor(a_1+\dots+a_r)/2\rfloor$ prime factors. (Start with a grid
$\{0,\dots,p\}\times\{0,\dots,q\}$.)
\end{exercise}

\section{Hints}

\begin{hint}{orders:ex:klein}
Look at the convex hull of the five points. The hard case is a triangle with two points inside.
\end{hint}
\begin{hint}{orders:ex:esproduct}
Mirsky's theorem, with the order of the proof of Theorem~\ref{orders:thm:es}.
\end{hint}
\begin{hint}{orders:ex:ackermannvalues}
Induction on $n$ in each line, using the line above. $2^{65536}$ has $\lfloor65536\log_{10}2\rfloor+1$ digits.
\end{hint}
\begin{hint}{orders:ex:rewriting}
Write the word as $a^{c_0}ba^{c_1}b\cdots ba^{c_k}$; a step at the $j$-th $b$ changes $c_j$ and $c_{j-1}$. Order the vectors $(c_k,\dots,c_1)$
lexicographically.
\end{hint}
\begin{hint}{orders:ex:hydra}
Encode a tree as the multiset of the encodings of its subtrees at the children of the root, and use
Theorem~\ref{orders:thm:multiset} once per level. The depth of the hydra never increases.
\end{hint}
\begin{hint}{orders:ex:fixedpoints}
(b) Try a map that jumps at $\frac12$.
\end{hint}
\begin{hint}{orders:ex:youngtableaux}
Write a linear extension as the word listing, at each step, the row of the next element.
\end{hint}
\begin{hint}{orders:ex:dedekind}
To a monotone function, associate the minimal points where it is $1$.
\end{hint}
\begin{hint}{orders:ex:hansel}
A monotone function restricted to a saturated chain is $0$ up to some point and $1$ after it.
\end{hint}
\begin{hint}{orders:ex:divisorscd}
In the grid, use hooks: a vertical segment followed by a horizontal one. Then show that the product of a poset having a
symmetric chain decomposition with a chain has one too.
\end{hint}

\section{Solutions}

\begin{solution}{orders:ex:klein}
\idea Four points are in convex position when each of them is a vertex of their convex hull, and a point is a hull vertex
when some line through it leaves the other points on one side. The convex hull of the five points tells us where to look.
\writeup If the convex hull of the five points has $5$ or $4$ vertices, any four of them, or these four, form a convex
quadrilateral. Otherwise the hull is a triangle $ABC$ with two points $D,E$ inside. The line $DE$ passes through no vertex, as no three
points are collinear, so two of the vertices, say $A$ and $B$, lie strictly on the same side of it. Then the line $DE$ leaves $A$ and $B$
on one side, so $D$ and $E$ are vertices of the convex hull of $\{A,B,D,E\}$; and the line $AB$ leaves $D$ and $E$ on one side (the side of
the triangle), so $A$ and $B$ are vertices of that hull too. All four points are hull vertices: $A,B,D,E$ form a convex
quadrilateral.
\end{solution}

\begin{solution}{orders:ex:esproduct}
\idea Mirsky covers the sequence by as many antichains as the longest chain is long, and each antichain is short.
\writeup With the order of Theorem~\ref{orders:thm:es}, chains are increasing and antichains decreasing subsequences. If the longest
increasing subsequence has length $h$ and the longest decreasing one length $d$, Mirsky's theorem partitions the $N$ terms into $h$
antichains of size at most $d$, so $N\le hd$. The blocks example of Theorem~\ref{orders:thm:es} shows equality is possible when $N=hd$.
\end{solution}

\begin{solution}{orders:ex:ackermannvalues}
\idea Each row of Ackermann's function iterates the previous one: $A(m+1,n)$ is $A(m,\cdot)$ applied $n+1$ times to $1$.
\writeup Induction on $n$. $A(1,0)=A(0,1)=2$ and $A(1,n+1)=A(0,A(1,n))=A(1,n)+1$, so $A(1,n)=n+2$. $A(2,0)=A(1,1)=3$ and
$A(2,n+1)=A(2,n)+2$, so $A(2,n)=2n+3$. $A(3,0)=A(2,1)=5=2^3-3$ and $A(3,n+1)=2A(3,n)+3$, so $A(3,n)+3$ doubles at each step:
$A(3,n)=2^{n+3}-3$. $A(4,0)=A(3,1)=13=2^{2^2}-3$ and $A(4,n+1)+3=2^{A(4,n)+3}$, so $A(4,n)+3$ is a tower of $n+3$ twos. Then
$A(4,1)=2^{16}-3=65533$ and $A(4,2)=2^{65536}-3$, which has $\lfloor65536\log_{10}2\rfloor+1=19\,729$ digits ($2^{65536}$ does not end in a string of
zeros, so subtracting $3$ keeps the number of digits).
\end{solution}

\begin{solution}{orders:ex:rewriting}
\idea The word gets longer, so length is useless. But a step moves an $a$ leftwards past a $b$: the $a$'s drift to the left,
and the counts of $a$'s between the $b$'s, read from the right, decrease lexicographically.
\writeup The number $k$ of $b$'s never changes. Write the word as $a^{c_0}ba^{c_1}b\cdots ba^{c_k}$. Replacing $ba$ by $aab$ at the $j$-th $b$ (which needs
$c_j\ge1$) turns $c_j$ into $c_j-1$ and $c_{j-1}$ into $c_{j-1}+2$, other counts unchanged. So the vector $(c_k,c_{k-1},\dots,c_1)$ decreases in the
lexicographic order of $\N^k$ (the first changed coordinate, $c_j$, decreases; $c_0$ is not recorded), which is well-founded by
Proposition~\ref{orders:prop:lex} and induction. Hence the process stops, and it stops exactly when $c_1=\dots=c_k=0$, at $a^{c_0}b^k$.
From $b^ka$, an $a$ sitting after $j\ge1$ letters $b$ is eventually moved once and becomes two $a$'s after $j-1$ letters $b$; so the
final word is $a^{2^k}b^k$, and the number of steps is the number of moves, $1+2+\dots+2^{k-1}=2^k-1$, whatever the order.
\end{solution}

\begin{solution}{orders:ex:hydra}
\idea The hydra grows, but it only grows \emph{lower} copies of what was cut: in the right order, cutting always makes it
smaller. Represent a node by the multiset of the representations of its children; then a cut replaces one element, at each
level from the cut up to the root, by smaller ones.
\writeup The cut and the regrowth never increase the depth of the tree, since the copies of the subtree at $p$ hang at $p$'s own
depth. So fix $D$, the initial depth. Define sets $M_0\subseteq M_1\subseteq\cdots$ and orders: $M_0=\{\emptyset\}$, and $M_{d+1}$ is the set of finite multisets of
elements of $M_d$, with the multiset order $<_{d+1}$ induced by $<_d$. By Theorem~\ref{orders:thm:multiset} and induction on $d$, no $<_d$ has an
infinite decreasing sequence. Encode a node $v$ whose subtree has depth at most $d$ by $\varphi(v)\in M_d$: a leaf is $\emptyset$, and an
inner node is the multiset of the codes of its children. (One checks by induction that this is consistent across the
$M_d$, since $M_d\subseteq M_{d+1}$ and the order $<_{d+1}$ restricted to $M_d$ is $<_d$.)

Consider a step. At $p$, the head $h$ (code $\emptyset$) is removed and nothing added, so the new code $\varphi'(p)$ is smaller than $\varphi(p)$. If $p$ is the
root, we are done. Otherwise, at $g$ the element $\varphi(p)$ is removed and $n+1$ copies of $\varphi'(p)$, each smaller, are added: $\varphi(g)$
decreases. At each ancestor of $g$, one element (the code of a child on the path) is replaced by a smaller one: the code
decreases. So the code of the root decreases at every step in $(M_D,<_D)$, which has no infinite decreasing sequence. The
fight ends, and since Hercules can always cut while a head exists, it ends with the bare root.
\end{solution}

\begin{remark}
Kirby and Paris proved that this statement, which only speaks of finite trees, cannot be proved in Peano arithmetic: any
proof must use an argument equivalent to the well-foundedness of the ordinal $\varepsilon_0$, which is what the nested multisets
encode. And the fight is long: for a hydra of depth $4$ it can outlast any number that you can write down. This is an example of a finite combinatorial statement that formal arithmetic cannot prove.
\end{remark}

\begin{solution}{orders:ex:fixedpoints}
\idea Without monotonicity a map can simply swap the bottom and the top. With monotonicity a fixed point exists, but the
iterates from $0$ may creep towards it forever, or even towards a point that is not fixed, if $f$ jumps.
\writeup (a) $f(0)=1$ and $f(1)=0$. (b) Let $f(x)=\frac{1+x}2$. It is increasing, its only fixed point is $1$, and $f^k(0)=1-2^{-k}$ never reaches
it; here the limit of the iterates is the least fixed point. It need not be: let $f(x)=\frac x2+\frac14$ for $x<\frac12$ and $f(x)=1$ for
$x\ge\frac12$. This $f$ is non-decreasing, since $f(x)<\frac12$ for $x<\frac12$. A fixed point $x<\frac12$ would satisfy $\frac x2+\frac14=x$, that is, $x=\frac12$, which is
absurd; so the only fixed point is $1$. The iterates $f^k(0)=\frac12-2^{-k-1}$ increase to $\frac12$, which is not a fixed point, and one more
step, $f(\frac12)=1$, reaches the least fixed point. In general the iteration must be continued ``beyond infinity'', through the
ordinals; when $f$ preserves limits of increasing sequences (Scott continuity) the limit is the least fixed point, and
this is the setting of the semantics of recursive programs.
\end{solution}

\begin{solution}{orders:ex:youngtableaux}
\idea A linear extension lists the $2n$ cells one at a time; reading only the rows, it is a word of $n$ ones and $n$ twos, and
the constraints say a two can never overtake the ones.
\writeup In a linear extension, the elements of each row appear in increasing order, so the extension is determined by the word
$w_1\cdots w_{2n}$ giving the row of each successive element. Such a word has $n$ letters $1$ and $n$ letters $2$, and the constraint
$(1,j)<(2,j)$ says that the $j$-th $2$ comes after the $j$-th $1$, that is, every prefix has at least as many $1$s as $2$s. Conversely
every such word gives a linear extension. Replacing $1$ by an up-step and $2$ by a down-step, these words are the Dyck paths of
length $2n$, counted by the Catalan number $\frac1{n+1}\binom{2n}n$ (Volume~3, Chapter~3).
\end{solution}

\begin{solution}{orders:ex:dedekind}
\idea A monotone function is determined by its minimal true points (Exercise~\ref{orders:ex:monotone}), and minimal points form an antichain.
\writeup Identify $\{0,1\}^n$ with $B_n$. To a monotone $f$ associate the set $\mathcal A_f$ of minimal sets $A$ with $f(A)=1$, an antichain; conversely an
antichain $\mathcal A$ defines the monotone function ``$f(X)=1$ if $X$ contains a member of $\mathcal A$''. These maps are inverse to each other (the
empty antichain corresponds to $f\equiv0$, the antichain $\{\emptyset\}$ to $f\equiv1$). For $n=2$: $\emptyset$, $\{\emptyset\}$, $\{\{1\}\}$, $\{\{2\}\}$, $\{\{1\},\{2\}\}$,
$\{\{1,2\}\}$: six. For $n=3$, count by the size of the smallest member: the two antichains $\emptyset$ and $\{\emptyset\}$; those whose smallest
sets are singletons; and so on. The total is $20$, and the next values are $168$, $7581$, $7\,828\,354$: the \emph{Dedekind numbers}, known
only up to $n=9$ (computed in 2023).
\end{solution}

\begin{solution}{orders:ex:hansel}
\idea Antichains of the middle rank give the lower bound; for the upper bound, a symmetric chain decomposition cuts $B_n$ into
$\binom n{\lfloor n/2\rfloor}$ chains on each of which a monotone function is a single threshold.
\writeup Every subset of the middle rank is an antichain: $2^{\binom n{\lfloor n/2\rfloor}}$ of them. By Exercise~\ref{orders:ex:dedekind}, $D(n)$ counts monotone functions. A
monotone function is determined by its restrictions to the chains of a symmetric chain decomposition
(Theorem~\ref{orders:thm:scd}). On a chain $A_1\subset\dots\subset A_m$ a monotone function is $0$ up to some index and $1$ after it: $m+1\le n+2$ possibilities. With
$\binom n{\lfloor n/2\rfloor}$ chains, $D(n)\le(n+2)^{\binom n{\lfloor n/2\rfloor}}$. Hansel (1966) refined the argument, using that the values on shorter chains leave at most
three possibilities on each chain, to get $D(n)\le3^{\binom n{\lfloor n/2\rfloor}}$, and Korshunov (1981) found the asymptotics of $D(n)$.
\end{solution}

\begin{solution}{orders:ex:divisorscd}
\idea A grid is a product of two chains; it splits into ``hooks'' which are symmetric. A product of a symmetric chain with a chain
is a grid, so symmetric decompositions survive multiplication by chains, and the divisor lattice is a product of chains.
\writeup \emph{Grids.} In $G=\{0,\dots,p\}\times\{0,\dots,q\}$ with rank $x+y$, for $0\le j\le\min(p,q)$ let $H_j$ be the vertical segment from $(j,0)$ to $(j,q-j)$
followed by the horizontal segment from $(j+1,q-j)$ to $(p,q-j)$. It is saturated, from rank $j$ to rank $p+q-j$, hence symmetric. A
point $(x,y)$ lies on the vertical part of $H_x$ if $y\le q-x$, and otherwise on the horizontal part of $H_{q-y}$ (then $0\le q-y<x\le p$); it
lies on no other hook. So the hooks partition $G$.

\emph{Products.} Let $P$ be a ranked poset of rank $R$ with a symmetric chain decomposition, and $L=\{0,\dots,\ell\}$ a chain. For a chain $C$
of $P$ from rank $r$ to rank $R-r$, $C\times L$ is a grid $\{0,\dots,R-2r\}\times\{0,\dots,\ell\}$ shifted by rank $r$. Its hooks run from rank $r+j$ to rank
$r+(R-2r)+\ell-j$, and these add up to $R+\ell$, the rank of $P\times L$: they are symmetric in $P\times L$, and saturated there. So $P\times L$ has a
symmetric chain decomposition.

\emph{Divisors.} Sending $d=\prod p_i^{b_i}$ to $(b_1,\dots,b_r)$ identifies the divisors of $n$ with the product of chains $\prod\{0,\dots,a_i\}$, rank being
$\sum b_i$. By induction on $r$, it has a symmetric chain decomposition. Each chain meets the middle rank $\lfloor(\sum a_i)/2\rfloor$ exactly
once, so the number of chains equals the size of that rank; an antichain meets each chain at most once, so no antichain is
larger. For $n=60$ this is Figure~\ref{orders:fig:divisors}: the largest antichain is $\{4,6,10,15\}$.
\end{solution}

\section*{Sources and further reading}

\begin{itemize}
\item P.~Erd\H os and G.~Szekeres, \emph{A combinatorial problem in geometry}, Compositio Mathematica 2 (1935), 463--470.
The happy ending paper.
\item A.~Suk, \emph{On the Erd\H os--Szekeres convex polygon problem}, Journal of the AMS 30 (2017), 1047--1053; G.~Szekeres
and L.~Peters, \emph{Computer solution to the 17-point Erd\H os--Szekeres problem}, ANZIAM Journal 48 (2006), 151--164.
\item R.~P.~Dilworth, \emph{A decomposition theorem for partially ordered sets}, Annals of Mathematics 51 (1950), 161--166;
M.~A.~Perles, \emph{A proof of Dilworth's decomposition theorem for partially ordered sets}, Israel Journal of Mathematics 1
(1963), 105--107; L.~Mirsky, \emph{A dual of Dilworth's decomposition theorem}, American Mathematical Monthly 78 (1971), 876--877.
\item A.~Tarski, \emph{A lattice-theoretical fixpoint theorem and its applications}, Pacific Journal of Mathematics 5 (1955),
285--309.
\item N.~Dershowitz and Z.~Manna, \emph{Proving termination with multiset orderings}, Communications of the ACM 22 (1979),
465--476.
\item L.~Kirby and J.~Paris, \emph{Accessible independence results for Peano arithmetic}, Bulletin of the London
Mathematical Society 14 (1982), 285--293. The Hydra.
\item E.~Sperner, \emph{Ein Satz \"uber Untermengen einer endlichen Menge}, Mathematische Zeitschrift 27 (1928), 544--548;
D.~Lubell, \emph{A short proof of Sperner's lemma}, Journal of Combinatorial Theory 1 (1966), 299.
\item N.~G.~de Bruijn, C.~van Ebbenhorst Tengbergen and D.~Kruyswijk, \emph{On the set of divisors of a number}, Nieuw
Archief voor Wiskunde 23 (1951), 191--193.
\item P.~Erd\H os, C.~Ko and R.~Rado, \emph{Intersection theorems for systems of finite sets}, Quarterly Journal of
Mathematics 12 (1961), 313--320; G.~O.~H.~Katona, \emph{A simple proof of the Erd\H os--Chao Ko--Rado theorem}, Journal of
Combinatorial Theory B 13 (1972), 183--184.
\item B.~A.~Davey and H.~A.~Priestley, \emph{Introduction to Lattices and Order}, 2nd ed., Cambridge University Press, 2002.
Lattices, fixed points and their use in computer science.
\item I.~Anderson, \emph{Combinatorics of Finite Sets}, Oxford University Press, 1987. Sperner theory and symmetric chains.
\end{itemize}

\section{Closure rules and infinite trees}
\begin{exercise}[$\star\star$]\label{orders:ex:monotone}
Show that a Boolean function is monotone if and only if it is computed by a formula using only $\land$ and $\lor$ (or is a
constant). Give such a formula for the \emph{majority} of three bits.
\end{exercise}

\begin{solution}{orders:ex:monotone}
\idea A monotone function is determined by its minimal ``true'' points: it is $1$ exactly above one of them, and ``above $a$'' is
a conjunction of variables.
\writeup Conjunctions and disjunctions preserve monotonicity, which gives one direction. Conversely let $f$ be monotone and not constant, and let $A$ be the set of
minimal points $a$ with $f(a)=1$ (for the coordinatewise order). Since $f(0,\dots,0)=0$ (otherwise $f\equiv1$ by monotonicity), every
$a\in A$ has a coordinate $1$, and $t_a=\bigwedge_{a_i=1}x_i$ is a formula without negation, true at $b$ if and only if $a\le b$. Let
$\varphi=\bigvee_{a\in A}t_a$. If $f(b)=1$, there is a minimal $a\le b$ with $f(a)=1$ (go down while possible), and $t_a$ is true at $b$. If $\varphi$ is
true at $b$, then $a\le b$ for some $a\in A$, and $f(b)\ge f(a)=1$. So $\varphi$ computes $f$. For majority, the minimal true points are
$110,101,011$, giving $(x_1\land x_2)\lor(x_1\land x_3)\lor(x_2\land x_3)$.
\end{solution}

\begin{exercise}[$\star\star$]\label{orders:ex:konig}
\emph{K\H onig's lemma.} A rooted tree in which every vertex has finitely many children and which has infinitely many
vertices has an infinite path from the root. Prove it directly.
\end{exercise}

\begin{solution}{orders:ex:konig}
\idea An infinite set split into finitely many parts has an infinite part: the pigeonhole principle, applied at each step
down the tree.
\writeup \emph{Directly.} Call a vertex \emph{rich} if it has infinitely many descendants. The root is rich. A rich vertex has finitely
many children, and its descendants are the children and their descendants, so at least one child is rich. Starting
from the root and always moving to a rich child gives an infinite path.

\end{solution}
